{\Large\bfseries Nearest sites mathematical overview}\par
\textbf{Setting.} Take finitely many distinct sites in the Euclidean plane. A point belongs to every nearest site, so cells include shared boundaries. Sites are fixed within a map, and distance is to their centers. The printed frame is a window unless a problem explicitly restricts the target to it.

\textbf{Two sites and many sites.} For distinct sites $a,b$, points at equal distance form their perpendicular bisector; the closed half-plane on $a$'s side consists of points at least as close to $a$ as to $b$. The cell of $a$ is the intersection of these half-planes over every other site. It is convex and contains $a$, but can be unbounded. A pairwise equality is part of a shared cell boundary only where no third site is nearer. Three-way and higher ties are allowed.

\textbf{Changing and prescribing cells.} Adding a site adds a constraint to each old cell, so old cells can shrink or stay the same but cannot grow. To realize a convex polygon as the full cell of a site strictly inside it, reflect that site across each supporting edge line and place a competitor at each image. Intersecting the corresponding half-planes gives exactly the polygon. A nonconvex region cannot be one cell. Each distinct straight edge of a full polygonal cell requires a different competitor, giving a lower bound in minimum-site tasks.

\textbf{Across the bands.} K--1 compares lengths, records all tied letters and investigates simple placements. Grades 2--3 builds whole regions and sees how a third site changes a two-site boundary. Grades 4--5 explains half-plane intersections, insertion, convexity and inverse constructions. Coordinates and squared-distance algebra are adult checks, not prerequisites.

\textbf{Experiment and proof.} String comparisons and folds reveal patterns and suggest conjectures. Exact symmetry, perpendicular-bisector arguments and intersections justify an entire boundary or region. A few sampled points cannot prove a whole region is correct. Keep the theorem map with adults during the initial search; the preparation, longer explanations and all 24 problem solutions follow.
